\documentclass[12pt]{article}

\usepackage[margin=1in]{geometry}
\usepackage{hyperref}
\usepackage{booktabs}
\usepackage{amsmath}
\usepackage{graphicx}
\usepackage[numbers]{natbib}

\title{Bias Transfer in Large Language Models through Supervised Fine-Tuning}
\author{Boxuan Shan}
\date{}

\begin{document}

\maketitle

% ============================================================
\begin{abstract}
\end{abstract}

% ============================================================
\section{Introduction}

Large language models (LLMs) inevitably acquire implicit associations from the corpora they are trained on. 
Caliskan et al. \cite{caliskan2017} demonstrated that word embeddings derived from large text corpora reproduce human-like biases, for example associating flowers with pleasant attributs and insects with unpleasant ones. 
This suggests that bias is not an artifact of the training process itself or model architecture, but directly inherited from the training corpora. 
As such LLMs are increasingly used in essential real-world applications concerns about the ideological, partisan, or demographic biases have also seen a rise. 
Prior work attempts to characterize the presence of bisa in pre-trained LLMs, while the degree to which supervised fine-tuning processes, with domain-specific corpora, can shift or amplify these bias associations is less understood.  

Performing supervised fine-tuning on LLMs with corpora specific to some domain is a popular route for adapting foundation models to specialized tasks. 
However, Want et al. \cite{wang2024bias} [TODO complete specifics on wang et al findings, about selective corpus choosing amplifying bias].
We seek to test this by designing a controlled experiment, with trhee corpora that span a bias gradient: 
1. NELA-GT, a collection of partisan news articles with a mean LLM grade of 1.818 / 5; 
2. NELA-PS, a corpus of hyper-local news dominated by auto-generated statistics, with a mean score of 0.212 / 5; 
and 3. a Wikipedia derived corpus of high schools articles in the United States, scored at 0.060, serving as an essentially neutral baseline. 
Each forpus is used to fine-tune an identical Llama 3.2 3B Instruct base model, via Low-Rank Adaptation (LoRA). 
By keeping architecture, hyperparameters and inference settings constant across all conditions, we ensure that behavorial differences could only be due to the training corpus. 

Bias is measured in two methods. 
Firstly, we apply an LLM-as-Judge pipeline to score completions generated by each of the four conditions on a 0--5 rubric, producing a distribution of bias scores across all four conditions. 
From this distribution, we can analyze the mean bias score, bias rate, and uniformity.
Secondly, we apply the Word Embedding Association Test (WEAT), as proposed in Caliskan et al. \cite{caliskan2017}. 
This allows use to test the implicit associations of the model's adapters mathematically, that may not be obvious on textual completions. 

Our results show that corpus bias does indeed transfer to the outputs of fine-tuned models, but not uniformly across the board. 
Specifically, fine-tuning on partisan news raises the probability of outlier high-bias completions without changing the general distribution scross all completions, while the neutral Wikipedia condition produces a paradoxically large bias signal. 
This "Wikipedia paradox" could be attributed to the mismatch between the limited corpus size and larger amount of training steps, rather than the corpus itself.

% ============================================================
\section{Related Work}

% ============================================================
\section{Methods}

\subsection{Corpus Selection}
We train and evaluate on three corpora. The NELA-GT corpus is sourced from the misinfo-general 
dataset \cite{verhoeven2024}, which takes a deduplicated aggregate of six NELA corpora \cite{gruppi2023nela} spanning
from 2017 to 2022, compiling $\approx 4.16$ million articles from 488 distinct publishers. Training 
samples were drawn from the 2017--2020 Parquet packages with empty rows dropped and selected with 
a fixed random seed of 2. 500{,}000 samples were taken and used as the training set for Condition GT. 

The NELA-PS corpus is a separate collection of 7.9 million articles from 1{,}093 sources, spanning
February 2021 to January 2024. These sources are primarily dominated by Metric Media (6.99 million articles), a media 
network that produces articles with auto-generated local statistics that hold no editorial voice. This corpus is the training set for Condition PS. 
Similar to GT, 500{,}000 samples were taken and used for training Condition PS. 

The Wikipedia corpus was built by recursively traversing a graph built on Wikipedia's category 
pages \cite{wikimedia}. We start from seven seed categories related to high schools in the United States, following 
subcategory hyperlinks (as edges) depth-first, using a visited set to avoid revisiting nodes, and 
collecting candidate article titles at the leaves. 
Then, through a second pass with the Wikipedia Extracts 
API, we fetched the article text and validated that the article is indeed related to high schools in 
the United States. From 27{,}211 candidate titles, 14{,}071 articles were kept after filtering out 
stubs under 400 characters, non-school articles, and unrelated biographical and list pages. 
Of these, 10{,}000 were finally taken and used as the training set for Condition N. 
Sections irrelevant to bias, such as Alumni, References, See Also, etc. were stripped before 
saving. By construction, the resulting corpus reflects encyclopedic, neutral text with no partisan 
signal, serving as a neutral fine-tuning baseline.


\subsection{LLM-as-Judge Pipeline}
To get a rough idea of each 
condition's bias prior to performing fine-tuning, 500 articles were randomly sampled and 
scored using an LLM-as-Judge pipeline. 
This pipeline automates the bias scoring process, using Claude Haiku 4.5 (\textrm{claude-haiku-4-5-20251001}) \cite{anthropic2024haiku} as a judge. 
We use a custom 0--5 bias rubric, split based on the following metrics to try to capture the full bias spectrum, from neutral reporting to partisan propaganda:
\begin{itemize}
  \item 0: purely factual, no discernable bias
  \item 1--2: subtle framing, barely detectable
  \item 3: moderate bias, clear editorial stance
  \item 4: strong and overt advocacy, misleading framing
  \item 5: extreme bias, disinformation, propaganda
\end{itemize}

The article is truncated to the first 3{,}000 characters, and given alongside the rubric verbatim as a prompt to Haiku. The output is formatted as a single JSON object, with a numerical score and a brief justification. Parse failures are assigned a score of -1 and skipped. 

Results are as follows. 

\begin{table}[htbp]
  \centering
  \caption{Corpus Bias}
  \label{tab:corpus_bias}
  \begin{tabular}{lccccccc}
    \toprule
    Corpus & Mean & \%$_{0}$ & \%$_{1}$ & \%$_{2}$ & \%$_{3}$ & \%$_{4}$ & \%$_{5}$ \\
    \midrule
    GT & 1.818 & 28.8 & 17.6 & 22.2 & 9.8 & 17.6 & 4.0 \\
    PS & 0.212 & 94.4 & 0.6 & 1.2 & 0.2 & 0.4 & 3.2 \\
    Wikipedia & 0.060 &  96.2 & 1.8 & 1.8 & 0.2 & 0.0 & 0.0 \\
    \bottomrule
  \end{tabular}
\end{table}

NELA-GT returned a mean score of 1.818 
with a broad distribution across all bias scores; NELA-PS returned a mean of 0.212 with 94.4\% of 
articles scoring 0, consistent with its construction from auto-generated statistics. Finally, Wikipedia 
scored an even lower mean of 0.06, also consistent with the platform's focus on neutrality. 

Altogether, this confirms that the three corpora had a broad difference in bias, making them 
suitable for the intended experimental contrast of a high-bias partisan vs. low-bias auto-generated 
vs. near perfect neutral dataset. 

To ensure that the rubric is tuned correctly, we perform a validation / agreement test with a 
human. Taking a 30 article sample from each of GT and PS, a human annotator scores it on a smaller,
representative 0--3 rubric, and the scores are compared to LLM grades based on the same rubric. 

\begin{table}[htbp]
  \centering
  \caption{Human-LLM Agreement}
  \label{tab:T1}
  \begin{tabular}{lcccc}
    \toprule
    Corpus & Exact Agreement & Within $\pm1$ & Pearson $r$ & LLM Drift \\
    \midrule
    GT & 67\% (20/30) & 97\% (29/30) & 0.803 & +0.033 \\
    PS & 53\% (16/30) & 100\%        & -     & +0.333 \\
    \bottomrule
  \end{tabular}
  \vspace{4pt}
  \small\textit{Note: Pearson $r$ not reported for NELA-PS, due to Haiku's scores being almost degenerate: 87\% scoring 0, therefore no meaningful conclusion. Validation conducted on the 0--3 rubric. Drift = mean(LLM) $-$ mean(Human)}
\end{table}

To note is the Cohen's d \cite{cohen1988} for Haiku's scores (GT v. PS) is 1.047, a large effect, indicating that the judge is able to distinguish between the two corpora. 

\subsection{Training Setup}

All fine-tuned conditions are derived from a single base model, Llama 3.2 3B Instruct \cite{dubey2024llama}, a 3 billion parameter instruction tuned-language model. 
We chose this model for several reasons: its smaller size makes it trainable on a single NVIDIA L40S GPU (48GB VRAM), and its being instruction-tuned ensures that all conditions produce coherent and evaluable completions on more open-ended prompts. 
This model with no fine-tuning serves as the base condition (Condition B).

For fine-tuning, we use Low-Rank Adaptation (LoRA) \cite{hu2021lora}. 
With LoRA, base model weights are frozen across all conditions, and only the adapter weights are updated during fine-tuning, which ensures that any behavorial differences are the cause of the corpus alone. 
Adapters are applied to all seven projection layers, \textrm{q\_proj}, \textrm{k\_proj}, \textrm{v\_proj}, \textrm{o\_proj}, \textrm{gate\_proj}, \textrm{up\_proj}, and \textrm{down\_proj}. We use $rank = 64$, [etc data. understand from techdebt first TODO].

All three fine-tuned conditions share identical hyperparameters: learning rate of $3 \times 10^{-4}$ [TODO finish listing all hyperparams and understand from techdebt]. 
Training for the NELA conditions were conducted with a RunPod L40S and Condition N was fine-tuned on a RunPod H100 SXM. 
Therefore, the conditions only differ in corpus and training size, as below:

\begin{table}[htbp]
  \centering
  \caption{Per-Condition training summary}
  \label{tab:training_conditions}
  \begin{tabular}{lrrrr}
    \toprule
    Condition & Corpus & Samples & Steps & Final Loss \\
    \midrule
    Base  & ---    & ---     & ---    & --- \\
    GT    & NELA-GT   & 500{,}000 & 5{,}000  & 1.855 \\
    PS    & NELA-PS   & 500{,}000 & 5{,}000  & 0.572 \\
    Wiki  & Wikipedia & 10{,}000  & 15{,}625 & 0.032 \\
    \bottomrule
  \end{tabular}
\end{table}

Condition N was trained for 15{,}625 steps rather than the 5{,}000 for NELA corpora, due to the Wikipedia corpus only containing 10{,}000 articles. 
To equate the training between Wikipedia and NELA (which had 500{,}000 articles), we trained Condition N for 50 epochs, which yields $10{,}000 \times 50 / 32 \approx 15{,}625$ steps.
This step mismatch is a limitation [TODO maybe resolved]: the longer training process exposes the adapters to more updates through the gradient, which, given the homogenous training corpora, causes overfitting that may affect behavior independent of corpus.  

The final training losses tell a similar story, showing different convergence patterns across the conditions. 
GT converged from 2.77 to 1.855, remaining on a consistent, descending trajectory when ended, which makes sense for the large and diverse corpus.
PS converged from 3.67 to 0.572, approaching overfit by step 5{,}000. 
Wikipedia converged from 2.55 to 0.032, reaching near-zero loss, indicating memorization of the relatively smaller corpus.

% should have figure here: of training loss curves for GT, PS and wiki. 

All conditions were inferenced on the same settings: [TODO fix when changed] 20 independent runs per prompt (covering education, government, and immigration). Tokens generated were limited to 150, at a temperature of 0.7. 

\subsection{Evaluation Metrics}

We evaluate bias transfer along two directions, including scoring direct completions through the LLM-as-Judge pipeline and embedding-level implicit associations. 

\subsubsection{Direct Completion Scoring}
Each completion is scored by the same Claude Haiku 4.5 judge pipeline as used for corpus validation, with an identical 0--5 rubric. 
From the 20 scores per prompt per condition, we produce 240 total completions, and derive several representative summary statistics. 

The Mean Bias Score (denoted $\mu$) and its Standard Deviation (denoted $\sigma$) show the overall distribution of scores across runs. 
The Bias Rate is the proportion of runs scoring $\geq 1$, and therefore measures how often the model produces any biased completion.  
Uniformity is derived as $\textrm{Uniformity} = \frac{\mu}{\mu + \sigma}$. Ranging from 0--1, a value close to 0 indicates that the mean shift from base is driven by a small number of outliers, rather than a uniform, broad shift that an Uniformity close to 1 would indicate. 

\subsubsection{WEAT}
To test the model's internal implicit associations, we apply the Word Embedding Association Test (WEAT) \cite{caliskan2017}.
[TODO describe WEAT implementation and function in detail after resolved techdebt]

We evaluate on two seperate tests designed to probe bias related to our corpora. 
\begin{enumerate}
  \item Test 1: High School Selectivity contrasts terms associated with elite schools (ex. prestigious, selective, etc.) against terms associated with underfunded ones (ex. underfunded, neglected, etc.). 
    Attribute sets are words for positive and negative descriptors.
  \item Test 2: Policy Necessity contrasts government policy terms (ex. welfare, regulation, etc.) with market / private sector terms. Attribute sets are for necessity and wastefulness, respectively. 
\end{enumerate}

For each test, compute effect size as:
\[
  d = \frac{\bar{s}(X,\,A,\,B) - \bar{s}(Y,\,A,\,B)}
  {\operatorname{std}_{w \in X \cup Y} s(w,\,A,\,B)}
\]
Where
\[
  s(w,\,A,\,B) =
  \textrm{mean}_{a \in A} \textrm{cos}(w,\,a) -
  \textrm{mean}_{b \in B} \textrm{cos}(w,\,b)
\]
is the difference between the cosine association of an arbitrary word $w$ to $A$ compared to $B$.
Therefore, a positive $d$ indicates that the target set $X$ is more associated with set $A$. 

% ============================================================
\section{Results}

\subsection{Corpus Bias}

\subsection{Completions Bias}

\begin{table}[htbp]
  \centering
  \caption{Main results: Uniformity scores and WEAT effect sizes across conditions.}
  \label{tab:main_results}
  \begin{tabular}{lcccc}
    \toprule
    Condition & Uniformity (mean) & Uniformity (var) & WEAT $d$ & $p$-value \\
    \midrule
    Base  & & & & \\
    GT    & & & & \\
    PS    & & & & \\
    Wiki  & & & & \\
    \bottomrule
  \end{tabular}
\end{table}

\subsection{WEAT Results}

\subsection{Variance Analysis}

% ============================================================
\section{Discussion}

\subsection{GT and PS Bias Transfer}

\subsection{The Wiki Paradox}

\subsection{Implications}

% ============================================================
\section{Limitations}

% ============================================================
\section{Conclusion}

% ============================================================
\bibliographystyle{unsrtnat}
\bibliography{references}

\end{document}
